\IfFileExists{stacks-project.cls}{%
\documentclass{stacks-project}
}{%
\documentclass{amsart}
}

% For dealing with references we use the comment environment
\usepackage{verbatim}
\newenvironment{reference}{\comment}{\endcomment}
%\newenvironment{reference}{}{}
\newenvironment{slogan}{\comment}{\endcomment}
\newenvironment{history}{\comment}{\endcomment}

% For commutative diagrams we use Xy-pic
\usepackage[all]{xy}

% We use 2cell for 2-commutative diagrams.
\xyoption{2cell}
\UseAllTwocells

% We use multicol for the list of chapters between chapters
\usepackage{multicol}

% This is generally recommended for better output
\usepackage{lmodern}
\usepackage[T1]{fontenc}

% For cross-file-references
\usepackage{xr-hyper}

% Package for hypertext links:
\usepackage{hyperref}

% For any local file, say "hello.tex" you want to link to please
% use \externaldocument[hello-]{hello}
\externaldocument[introduction-]{introduction}
\externaldocument[conventions-]{conventions}
\externaldocument[sets-]{sets}
\externaldocument[categories-]{categories}
\externaldocument[topology-]{topology}
\externaldocument[sheaves-]{sheaves}
\externaldocument[sites-]{sites}
\externaldocument[stacks-]{stacks}
\externaldocument[fields-]{fields}
\externaldocument[algebra-]{algebra}
\externaldocument[brauer-]{brauer}
\externaldocument[homology-]{homology}
\externaldocument[derived-]{derived}
\externaldocument[simplicial-]{simplicial}
\externaldocument[more-algebra-]{more-algebra}
\externaldocument[smoothing-]{smoothing}
\externaldocument[modules-]{modules}
\externaldocument[sites-modules-]{sites-modules}
\externaldocument[injectives-]{injectives}
\externaldocument[cohomology-]{cohomology}
\externaldocument[sites-cohomology-]{sites-cohomology}
\externaldocument[dga-]{dga}
\externaldocument[dpa-]{dpa}
\externaldocument[sdga-]{sdga}
\externaldocument[hypercovering-]{hypercovering}
\externaldocument[schemes-]{schemes}
\externaldocument[constructions-]{constructions}
\externaldocument[properties-]{properties}
\externaldocument[morphisms-]{morphisms}
\externaldocument[coherent-]{coherent}
\externaldocument[divisors-]{divisors}
\externaldocument[limits-]{limits}
\externaldocument[varieties-]{varieties}
\externaldocument[topologies-]{topologies}
\externaldocument[descent-]{descent}
\externaldocument[perfect-]{perfect}
\externaldocument[more-morphisms-]{more-morphisms}
\externaldocument[flat-]{flat}
\externaldocument[groupoids-]{groupoids}
\externaldocument[more-groupoids-]{more-groupoids}
\externaldocument[etale-]{etale}
\externaldocument[chow-]{chow}
\externaldocument[intersection-]{intersection}
\externaldocument[pic-]{pic}
\externaldocument[weil-]{weil}
\externaldocument[adequate-]{adequate}
\externaldocument[dualizing-]{dualizing}
\externaldocument[duality-]{duality}
\externaldocument[discriminant-]{discriminant}
\externaldocument[derham-]{derham}
\externaldocument[local-cohomology-]{local-cohomology}
\externaldocument[algebraization-]{algebraization}
\externaldocument[curves-]{curves}
\externaldocument[resolve-]{resolve}
\externaldocument[models-]{models}
\externaldocument[functors-]{functors}
\externaldocument[equiv-]{equiv}
\externaldocument[pione-]{pione}
\externaldocument[etale-cohomology-]{etale-cohomology}
\externaldocument[proetale-]{proetale}
\externaldocument[relative-cycles-]{relative-cycles}
\externaldocument[more-etale-]{more-etale}
\externaldocument[trace-]{trace}
\externaldocument[crystalline-]{crystalline}
\externaldocument[spaces-]{spaces}
\externaldocument[spaces-properties-]{spaces-properties}
\externaldocument[spaces-morphisms-]{spaces-morphisms}
\externaldocument[decent-spaces-]{decent-spaces}
\externaldocument[spaces-cohomology-]{spaces-cohomology}
\externaldocument[spaces-limits-]{spaces-limits}
\externaldocument[spaces-divisors-]{spaces-divisors}
\externaldocument[spaces-over-fields-]{spaces-over-fields}
\externaldocument[spaces-topologies-]{spaces-topologies}
\externaldocument[spaces-descent-]{spaces-descent}
\externaldocument[spaces-perfect-]{spaces-perfect}
\externaldocument[spaces-more-morphisms-]{spaces-more-morphisms}
\externaldocument[spaces-flat-]{spaces-flat}
\externaldocument[spaces-groupoids-]{spaces-groupoids}
\externaldocument[spaces-more-groupoids-]{spaces-more-groupoids}
\externaldocument[bootstrap-]{bootstrap}
\externaldocument[spaces-pushouts-]{spaces-pushouts}
\externaldocument[spaces-chow-]{spaces-chow}
\externaldocument[groupoids-quotients-]{groupoids-quotients}
\externaldocument[spaces-more-cohomology-]{spaces-more-cohomology}
\externaldocument[spaces-simplicial-]{spaces-simplicial}
\externaldocument[spaces-duality-]{spaces-duality}
\externaldocument[formal-spaces-]{formal-spaces}
\externaldocument[restricted-]{restricted}
\externaldocument[spaces-resolve-]{spaces-resolve}
\externaldocument[formal-defos-]{formal-defos}
\externaldocument[defos-]{defos}
\externaldocument[cotangent-]{cotangent}
\externaldocument[examples-defos-]{examples-defos}
\externaldocument[algebraic-]{algebraic}
\externaldocument[examples-stacks-]{examples-stacks}
\externaldocument[stacks-sheaves-]{stacks-sheaves}
\externaldocument[criteria-]{criteria}
\externaldocument[artin-]{artin}
\externaldocument[quot-]{quot}
\externaldocument[stacks-properties-]{stacks-properties}
\externaldocument[stacks-morphisms-]{stacks-morphisms}
\externaldocument[stacks-limits-]{stacks-limits}
\externaldocument[stacks-cohomology-]{stacks-cohomology}
\externaldocument[stacks-perfect-]{stacks-perfect}
\externaldocument[stacks-introduction-]{stacks-introduction}
\externaldocument[stacks-more-morphisms-]{stacks-more-morphisms}
\externaldocument[stacks-geometry-]{stacks-geometry}
\externaldocument[moduli-]{moduli}
\externaldocument[moduli-curves-]{moduli-curves}
\externaldocument[examples-]{examples}
\externaldocument[exercises-]{exercises}
\externaldocument[guide-]{guide}
\externaldocument[desirables-]{desirables}
\externaldocument[coding-]{coding}
\externaldocument[obsolete-]{obsolete}
\externaldocument[fdl-]{fdl}
\externaldocument[index-]{index}

% Theorem environments.
%
\theoremstyle{plain}
\newtheorem{theorem}[subsection]{Theorem}
\newtheorem{proposition}[subsection]{Proposition}
\newtheorem{lemma}[subsection]{Lemma}

\theoremstyle{definition}
\newtheorem{definition}[subsection]{Definition}
\newtheorem{example}[subsection]{Example}
\newtheorem{exercise}[subsection]{Exercise}
\newtheorem{situation}[subsection]{Situation}

\theoremstyle{remark}
\newtheorem{remark}[subsection]{Remark}
\newtheorem{remarks}[subsection]{Remarks}

\numberwithin{equation}{subsection}

% Macros
%
\def\lim{\mathop{\mathrm{lim}}\nolimits}
\def\colim{\mathop{\mathrm{colim}}\nolimits}
\def\Spec{\mathop{\mathrm{Spec}}}
\def\Hom{\mathop{\mathrm{Hom}}\nolimits}
\def\Ext{\mathop{\mathrm{Ext}}\nolimits}
\def\SheafHom{\mathop{\mathcal{H}\!\mathit{om}}\nolimits}
\def\SheafExt{\mathop{\mathcal{E}\!\mathit{xt}}\nolimits}
\def\Sch{\mathit{Sch}}
\def\Mor{\mathop{\mathrm{Mor}}\nolimits}
\def\Ob{\mathop{\mathrm{Ob}}\nolimits}
\def\Sh{\mathop{\mathit{Sh}}\nolimits}
\def\NL{\mathop{N\!L}\nolimits}
\def\CH{\mathop{\mathrm{CH}}\nolimits}
\def\proetale{{pro\text{-}\acute{e}tale}}
\def\etale{{\acute{e}tale}}
\def\QCoh{\mathit{QCoh}}
\def\Ker{\mathop{\mathrm{Ker}}}
\def\Im{\mathop{\mathrm{Im}}}
\def\Coker{\mathop{\mathrm{Coker}}}
\def\Coim{\mathop{\mathrm{Coim}}}

% Boxtimes
%
\DeclareMathSymbol{\boxtimes}{\mathbin}{AMSa}{"02}

%
% Macros for moduli stacks/spaces
%
\def\QCohstack{\mathcal{QC}\!\mathit{oh}}
\def\Cohstack{\mathcal{C}\!\mathit{oh}}
\def\Spacesstack{\mathcal{S}\!\mathit{paces}}
\def\Quotfunctor{\mathrm{Quot}}
\def\Hilbfunctor{\mathrm{Hilb}}
\def\Curvesstack{\mathcal{C}\!\mathit{urves}}
\def\Polarizedstack{\mathcal{P}\!\mathit{olarized}}
\def\Complexesstack{\mathcal{C}\!\mathit{omplexes}}
% \Pic is the operator that assigns to X its picard group, usage \Pic(X)
% \Picardstack_{X/B} denotes the Picard stack of X over B
% \Picardfunctor_{X/B} denotes the Picard functor of X over B
\def\Pic{\mathop{\mathrm{Pic}}\nolimits}
\def\Picardstack{\mathcal{P}\!\mathit{ic}}
\def\Picardfunctor{\mathrm{Pic}}
\def\Deformationcategory{\mathcal{D}\!\mathit{ef}}
\usepackage{cleveref}
\crefname{originalreference}{original reference}{original references}
\Crefname{originalreference}{Original reference}{Original references}
\makeatletter
\newlabel{algebra-lemma-directed-colimit-exact}{{00DB}{source}{}{sourcepointer-1}{}}
\newlabel{algebra-lemma-directed-colimit-exact@cref}{{[originalreference][0][]00DB}{[1][0][]source}{}{}{}}
\makeatother
\crefname{originalreference}{original reference}{original references}
\Crefname{originalreference}{Original reference}{Original references}
\makeatletter
\newlabel{categories-definition-essentially-constant-diagram}{{05PU}{source}{}{sourcepointer-2}{}}
\newlabel{categories-definition-essentially-constant-diagram@cref}{{[originalreference][0][]05PU}{[1][0][]source}{}{}{}}
\makeatother
\crefname{originalreference}{original reference}{original references}
\Crefname{originalreference}{Original reference}{Original references}
\makeatletter
\newlabel{categories-lemma-characterize-essentially-constant-ind}{{05PY}{source}{}{sourcepointer-3}{}}
\newlabel{categories-lemma-characterize-essentially-constant-ind@cref}{{[originalreference][0][]05PY}{[1][0][]source}{}{}{}}
\makeatother
\crefname{originalreference}{original reference}{original references}
\Crefname{originalreference}{Original reference}{Original references}
\makeatletter
\newlabel{categories-lemma-cofinal-essentially-constant}{{0A1S}{source}{}{sourcepointer-4}{}}
\newlabel{categories-lemma-cofinal-essentially-constant@cref}{{[originalreference][0][]0A1S}{[1][0][]source}{}{}{}}
\makeatother
\crefname{originalreference}{original reference}{original references}
\Crefname{originalreference}{Original reference}{Original references}
\makeatletter
\newlabel{homology-lemma-five-lemma}{{05QB}{source}{}{sourcepointer-5}{}}
\newlabel{homology-lemma-five-lemma@cref}{{[originalreference][0][]05QB}{[1][0][]source}{}{}{}}
\makeatother
\crefname{originallemma}{original lemma}{original lemmas}
\Crefname{originallemma}{Original lemma}{Original lemmas}
\makeatletter
\newlabel{lemma-derived-functor}{{05SA}{source}{}{sourcepointer-6}{}}
\newlabel{lemma-derived-functor@cref}{{[originallemma][0][]05SA}{[1][0][]source}{}{}{}}
\makeatother
\crefname{originallemma}{original lemma}{original lemmas}
\Crefname{originallemma}{Original lemma}{Original lemmas}
\makeatletter
\newlabel{lemma-derived-inverts}{{05SB}{source}{}{sourcepointer-7}{}}
\newlabel{lemma-derived-inverts@cref}{{[originallemma][0][]05SB}{[1][0][]source}{}{}{}}
\makeatother
\crefname{originallemma}{original lemma}{original lemmas}
\Crefname{originallemma}{Original lemma}{Original lemmas}
\makeatletter
\newlabel{lemma-limit-triangles}{{05R9}{source}{}{sourcepointer-8}{}}
\newlabel{lemma-limit-triangles@cref}{{[originallemma][0][]05R9}{[1][0][]source}{}{}{}}
\makeatother
\crefname{originalreference}{original reference}{original references}
\Crefname{originalreference}{Original reference}{Original references}
\makeatletter
\newlabel{situation-derived-functor}{{05S8}{source}{}{sourcepointer-9}{}}
\newlabel{situation-derived-functor@cref}{{[originalreference][0][]05S8}{[1][0][]source}{}{}{}}
\makeatother
\title{RF defined at two vertices of a distinguished triangle is defined at the third}
\author{The Stacks Project Authors}
\date{stacks-project:derived.tex:local-snapshot-20261012; GNU FDL 1.2 or later}
\begin{document}
\maketitle
\noindent\textit{Editorial note.} Counted span is the author's own proof environment, source lines 4908-5016, following the statement at 4895-4906.\par\medskip
\noindent\textit{Editorial note.} The statement's closing 'Similarly for LF' is not executed in the proof: mirror of the RF argument (left derived functor via the opposite triangulated category), recorded as a mirror-level omission.\par\medskip
\noindent\textit{Allowed standard prerequisites (editorial).}\begin{itemize}
\item Induced morphism RF(f). lemma-derived-functor, cited at proof line 4911; complete proof at source lines (queued, pb 746), same chapter. Supplies the induced morphism RF(f) : A -> B between chosen values.
\item Reduction lemma for definedness. lemma-derived-inverts, cited at proof line 4965; short source lemma (pb 32) in the same chapter. Supplies that it suffices to prove RF defined at Z' with the arrow gamma inducing the isomorphism C -> RF(Z').
\item Category of triangles over a fixed triangle. lemma-limit-triangles, cited at proof lines 4975 and 4986; its span 1812-1952 partially overlaps published range 316 and its proof in the chapter is complete (cofinality statement included). Supplies the category I of triangles over (X', Y', Z') and the cofinality of I -> Z'/S.
\item Exactness of filtered colimits. algebra-lemma-directed-colimit-exact, cross-chapter citation at proof line 5005 (tags mechanism). Supplies that the left column of the Mor diagram is an exact sequence.
\item Five lemma. homology-lemma-five-lemma, cross-chapter citation at proof line 5007 (tags mechanism). Supplies that colim Mor(W, F(Z'')) -> Mor(W, C) is bijective.
\item Characterization of essentially constant ind-objects. categories-lemma-characterize-essentially-constant-ind part (4), cross-chapter citation at proof line 5014 (tags mechanism). Supplies that the bijection for all W means F(Z'') is essentially constant over I with value C.
\end{itemize}\medskip
\ifdefined\hypertarget\hypertarget{sourcepointer-1}{}\fi
\noindent\textit{Source reference.} Original reference 00DB: focused-seeds/geometry/S9-der-2-out-of-3-defined/ref-tags.txt, source lines 5--5. Navigation pointer only: the counted proof cites this statement by \textbackslash{}ref as a standing input; the deduction the counted proof has to make is carried out inside the counted span, and no part of that deduction is delegated to this pointer.\par\medskip
\ifdefined\hypertarget\hypertarget{sourcepointer-2}{}\fi
\noindent\textit{Source reference.} Original reference 05PU: focused-seeds/geometry/S9-der-2-out-of-3-defined/ref-tags.txt, source lines 6--6. Navigation pointer only: the counted proof cites this statement by \textbackslash{}ref as a standing input; the deduction the counted proof has to make is carried out inside the counted span, and no part of that deduction is delegated to this pointer.\par\medskip
\ifdefined\hypertarget\hypertarget{sourcepointer-3}{}\fi
\noindent\textit{Source reference.} Original reference 05PY: focused-seeds/geometry/S9-der-2-out-of-3-defined/ref-tags.txt, source lines 7--7. Navigation pointer only: the counted proof cites this statement by \textbackslash{}ref as a standing input; the deduction the counted proof has to make is carried out inside the counted span, and no part of that deduction is delegated to this pointer.\par\medskip
\ifdefined\hypertarget\hypertarget{sourcepointer-4}{}\fi
\noindent\textit{Source reference.} Original reference 0A1S: focused-seeds/geometry/S9-der-2-out-of-3-defined/ref-tags.txt, source lines 8--8. Navigation pointer only: the counted proof cites this statement by \textbackslash{}ref as a standing input; the deduction the counted proof has to make is carried out inside the counted span, and no part of that deduction is delegated to this pointer.\par\medskip
\ifdefined\hypertarget\hypertarget{sourcepointer-5}{}\fi
\noindent\textit{Source reference.} Original reference 05QB: focused-seeds/geometry/S9-der-2-out-of-3-defined/ref-tags.txt, source lines 9--9. Navigation pointer only: the counted proof cites this statement by \textbackslash{}ref as a standing input; the deduction the counted proof has to make is carried out inside the counted span, and no part of that deduction is delegated to this pointer.\par\medskip
\ifdefined\hypertarget\hypertarget{sourcepointer-6}{}\fi
\noindent\textit{Source reference.} Original lemma 05SA: sources/stacks-derived/pinned-original-source.tex, source lines 4789--4789. Navigation pointer only: the counted proof cites this statement by \textbackslash{}ref as a standing input; the deduction the counted proof has to make is carried out inside the counted span, and no part of that deduction is delegated to this pointer.\par\medskip
\ifdefined\hypertarget\hypertarget{sourcepointer-7}{}\fi
\noindent\textit{Source reference.} Original lemma 05SB: sources/stacks-derived/pinned-original-source.tex, source lines 4855--4855. Navigation pointer only: the counted proof cites this statement by \textbackslash{}ref as a standing input; the deduction the counted proof has to make is carried out inside the counted span, and no part of that deduction is delegated to this pointer.\par\medskip
\ifdefined\hypertarget\hypertarget{sourcepointer-8}{}\fi
\noindent\textit{Source reference.} Original lemma 05R9: sources/stacks-derived/pinned-original-source.tex, source lines 1813--1813. Navigation pointer only: the counted proof cites this statement by \textbackslash{}ref as a standing input; the deduction the counted proof has to make is carried out inside the counted span, and no part of that deduction is delegated to this pointer.\par\medskip
\ifdefined\hypertarget\hypertarget{sourcepointer-9}{}\fi
\noindent\textit{Source reference.} Original reference 05S8: sources/stacks-derived/pinned-original-source.tex, source lines 4736--4736. Navigation pointer only: the counted proof cites this statement by \textbackslash{}ref as a standing input; the deduction the counted proof has to make is carried out inside the counted span, and no part of that deduction is delegated to this pointer.\par\medskip
\begin{lemma}
\label{lemma-2-out-of-3-defined}
Assumptions and notation as in
Situation \ref{situation-derived-functor}.
Let $(X, Y, Z, f, g, h)$ be a distinguished triangle of $\mathcal{D}$.
If $RF$ is defined at two out of three of $X, Y, Z$, then it is defined
at the third. Moreover, in this case
$$
(RF(X), RF(Y), RF(Z), RF(f), RF(g), RF(h))
$$
is a distinguished triangle in $\mathcal{D}'$. Similarly for $LF$.
\end{lemma}

\begin{proof}
Say $RF$ is defined at $X, Y$ with values $A, B$.
Let $RF(f) : A \to B$ be the induced morphism, see
Lemma \ref{lemma-derived-functor}.
We may choose a distinguished triangle
$(A, B, C, RF(f), b, c)$
in $\mathcal{D}'$. We claim that $C$ is a value of $RF$ at $Z$.

\medskip\noindent
To see this pick $s : X \to X'$ in $S$ such that there exists a morphism
$\alpha : A \to F(X')$ as in
Categories,
Definition \ref{categories-definition-essentially-constant-diagram}.
We may choose a commutative diagram
$$
\xymatrix{
X \ar[d]_f \ar[r]_s & X' \ar[d]^{f'} \\
Y \ar[r]^{s'} & Y'
}
$$
with $s' \in S$ by MS2. Using that $Y/S$ is filtered we can (after replacing
$s'$ by some $s'' : Y \to Y''$ in $S$) assume that there exists
a morphism $\beta : B \to F(Y')$ as in
Categories,
Definition \ref{categories-definition-essentially-constant-diagram}.
Picture
$$
\xymatrix{
A \ar[d]_{RF(f)} \ar[r]_-\alpha &
F(X') \ar[r] \ar[d]^{F(f')} &
A \ar[d]^{RF(f)} \\
B \ar[r]^-\beta & F(Y') \ar[r] & B
}
$$
It may not be true that the left square commutes, but the outer and
right squares commute.
The assumption that the ind-object $\{F(Y')\}_{s' : Y' \to Y}$
is essentially constant means that there exists a $s'' : Y \to Y''$
in $S$ and a morphism $h : Y' \to Y''$ such that $s'' = h \circ s'$ and
such that $F(h)$ equal to $F(Y') \to B \to F(Y') \to F(Y'')$. Hence
after replacing $Y'$ by $Y''$ and $\beta$ by $F(h) \circ \beta$ the
diagram will commute (by direct computation with arrows).

\medskip\noindent
Using MS6 choose a morphism of triangles
$$
(s, s', s'') : (X, Y, Z, f, g, h) \longrightarrow (X', Y', Z', f', g', h')
$$
with $s'' \in S$. By TR3 choose a morphism of triangles
$$
(\alpha, \beta, \gamma) :
(A, B, C, RF(f), b, c)
\longrightarrow
(F(X'), F(Y'), F(Z'), F(f'), F(g'), F(h'))
$$

\medskip\noindent
By Lemma \ref{lemma-derived-inverts} it suffices to prove that $RF(Z')$
is defined and that the arrow $\gamma : C \to F(Z')$ induces an isomorphism
$C \to RF(Z')$. Namely, then we will get an isomorphism
$$
(A, B, C, RF(f), b, c)
\longrightarrow
(RF(X'), RF(Y'), RF(Z'), RF(f'), RF(g'), RF(h'))
$$
of triangles and by TR1 we conclude that the target is a
distinguished triangle. Consider the category $\mathcal{I}$ of
Lemma \ref{lemma-limit-triangles}
of triangles
$$
\mathcal{I} =
\{(t, t', t'') : (X', Y', Z', f', g', h') \to (X'', Y'', Z'', f'', g'', h'')
\mid (t, t', t'') \in S\}
$$
To show that the system $F(Z'')$ is essentially constant over the category
$Z'/S$ is equivalent to showing that the system of $F(Z'')$ is essentially
constant over $\mathcal{I}$ because $\mathcal{I} \to Z'/S$ is cofinal, see
Categories, Lemma \ref{categories-lemma-cofinal-essentially-constant}
(cofinality is proven in Lemma \ref{lemma-limit-triangles}).
For any object $W$ in $\mathcal{D}'$ we
consider the diagram
$$
\xymatrix{
\colim_\mathcal{I} \Mor_{\mathcal{D}'}(W, F(X'')) \ar[d] &
\Mor_{\mathcal{D}'}(W, A) \ar[l] \ar[d] \\
\colim_\mathcal{I} \Mor_{\mathcal{D}'}(W, F(Y'')) \ar[d] &
\Mor_{\mathcal{D}'}(W, B) \ar[d] \ar[l] \\
\colim_\mathcal{I} \Mor_{\mathcal{D}'}(W, F(Z'')) \ar[d] &
\Mor_{\mathcal{D}'}(W, C) \ar[d] \ar[l] \\
\colim_\mathcal{I} \Mor_{\mathcal{D}'}(W, F(X''[1])) \ar[d] &
\Mor_{\mathcal{D}'}(W, A[1]) \ar[d] \ar[l] \\
\colim_\mathcal{I} \Mor_{\mathcal{D}'}(W, F(Y''[1])) &
\Mor_{\mathcal{D}'}(W, B[1]) \ar[l]
}
$$
where the horizontal arrows are given by composing with
$(\alpha, \beta, \gamma)$. Since filtered colimits are exact
(Algebra, Lemma \ref{algebra-lemma-directed-colimit-exact}) the left column
is an exact sequence. Thus the $5$ lemma
(Homology, Lemma \ref{homology-lemma-five-lemma}) tells us the map
$$
\colim_\mathcal{I} \Mor_{\mathcal{D}'}(W, F(Z''))
\longrightarrow
\Mor_{\mathcal{D}'}(W, C)
$$
is bijective. We conclude that $F(Z'')$ is essentially constant
over $\mathcal{I}$ with value $C$ by part (4) of
Categories, Lemma \ref{categories-lemma-characterize-essentially-constant-ind}.
\end{proof}

\end{document}
